\documentclass[12pt]{article}
\usepackage{styles/cwilliams-standard}

% Set up document metadata
\setauthors{Tzu-Chen Chiu, Arjun Sivakumar, Mara Schimmel, Christopher Williams}
\setheaderauthors{Chiu, Sivakumar, Schimmel, Williams}
\settitle{Sprint 4, Semester 2 Report}
\setclass{IPD}

\begin{document}

\maketitlepage

\section{Original Sprint 4 Backlog}

\begin{figure}[H]
    \centering
    \includegraphics[width=1\textwidth]{images/sprint4_backlog.png}
    \caption{Spring 3 Backlog}
\end{figure}

\section{Sprint 4 Backlog - End of Sprint Status}

\begin{figure}[H]
    \centering
    \includegraphics[width=1\textwidth]{images/sprint4_backlog_end.png}
    \caption{Spring 4 Backlog Updated}
\end{figure}

\section{Project Management Screenshots}

The Sprint backlogs in Google Sheets are being used to manage tasks.

\section{Sprint Retrospective}

\subsection*{Mara}

I worked on using the ARC, managing to run a basic python code. I also created two versions of
 data display using Streamlit, modeling possible layouts and UI for the final project. Next sprint 
 I plan on doing more examples on the ARC, and working more on the website.

\subsection*{Chris}

I continued my work from last week to revisit and understand with a greater depth MWPM 
and the data that we are embedding to ensure both the data we are using is correct, and that
 the model can perform as best and as authentically as possible in regards to AlphaQubit. 
 I also created a chart showcasing MWPM’s Logical Error Rate across different physical qubit 
 error rates at different code distances to showcase to our mentor to ensure my own 
 and our understanding as a group.

Next, we will start utilizing Google’s published data to do the same test before beginning to 
do more model trainings and attempt to replicate Google’s training process and tweak our 
model. I will revisit the architecture of the model with Arjun and Tzu-Chen and continue 
to provide data to Mara so that she can create a visualization portfolio of our project.


\subsection*{Arjun}

During this sprint, I focused on studying the Google dataset and thoroughly analyzing the AlphaQubit
 papers to gain a deeper understanding of the transformer architecture, parameter choices, training 
 setup, and implementation details. I was able to complete the tasks outlined in the Gantt chart for 
 this phase, while also supporting teammates by helping set up ARC workflows and contributing to 
 preparation for mentor meetings. A key impediment was the complexity of the dataset, particularly
  understanding the quantum circuits and underlying error model, which required structured mentor 
  meetings and team discussions, which clarified several conceptual gaps. No ethical concerns 
  arose, as we worked strictly with publicly available research data and maintained proper attribution.
   Overall, the sprint went well in terms of knowledge consolidation and team coordination, and
    the next phase will focus on fine-tuning the model using the exact parameters, curriculum design, 
    and implementation details derived from our analysis.

\subsection*{Tzu-Chen}

I completed a study of the Google dataset we’re using. This dataset is a distance-3 surface code 
with a single round. Each folder includes the circuit and DEM files, which allow us to generate 
measurement data for both data and ancilla qubits. From the measurement record, we can then 
compute detection events.

To ensure we understand the full pipeline, I reconstructed the measurement-to-event process
 step by step. I drew a diagram based on the circuit, identified the mapping between qubits and 
 measurement indices, and then filled the measurement outcomes into the corresponding rec[ ] 
 entries. Using the detector definitions in the circuit, I computed each detector outcome via XOR 
 and also incorporated the sweep terms. This took significant time to establish the relationship 
 between the raw data and the final detector events, and I’ve built an SOP for it.

Next, I’m starting to design and prepare the dataset splits for pretraining, validation, and testing. 
My plan is to use the DEM provided by Google as the base model, but modify the error-gate 
probabilities to generate additional training distributions. In particular, I want the dataset design 
to cover three key factors:

\begin{enumerate}
    \item code distance (d = 3 and 5)
    \item number of rounds (r = 1, 3, 5, ..., 25) and
    \item measurement basis (bX and bZ)
\end{enumerate}

\section{Product Backlog (Updated)}

\begin{figure}[H]
    \centering
    \includegraphics[width=1\textwidth]{images/sprint2_product_backlog_sem2_eos.png}
    \caption{Product Backlog - End of Sprint 2}
\end{figure}

\begin{aside}{Links}
For verification, \href{https://docs.google.com/spreadsheets/d/1E0FIjbewFiQXMTeIKEnVrVA3a1QH_6Hv9CjAknl78Ds/edit?gid=1121195217#gid=1121195217}{here} is the link to our Google Sheets. Under the ``Project Backlog'' is the image 
that is referenced above.
\end{aside}

\section{Burndown Chart}

\begin{figure}[H]
    \centering
    \includegraphics[width=1\textwidth]{images/burndown_chart_sprint2_sem2.png}
    \caption{Burndown Chart as of End of Sprint 2}
\end{figure}

\section{Issue Tracking}

\begin{table}[H]
\centering
\begin{tabular}{@{}llll@{}}
\toprule
\textbf{Issue} & \textbf{Description} & \textbf{Status} & \textbf{Date Comp.} \\ \midrule
ARC Access & Getting all members access to ARC & Completed & 02/16/2026 \\
ARC Bottlenecks & Setup, allocations, and constraints & Ongoing & Ongoing \\
Ensuring data quality & Ensure correct data & Ongoing & Ongoing \\
Creating better architecture & $\leftarrow$ & Complete & 2/5/2026 \\
Revising MWPM & Reviewing materials & Ongoing & Ongoing \\
Google Data Overview & Understanding Google's Data & Completed & 02/19/2026
\end{tabular}
\caption{Issue Tracking Summary}
\end{table}

\section{Sprint 3 Backlog - Semester 2}

\begin{figure}[H]
    \centering
    \includegraphics[width=1\textwidth]{images/sprint3_backlog_sem2.png}
    \caption{Sprint 3 Backlog}
\end{figure}

\end{document}